\section{Anexos}

\subsection{Reproducción del trabajo}
\label{anexo:reproduccion}

Todos los resultados del informe pueden regenerarse a partir de los archivos entregados.
Se entrega el comprimido \texttt{io2-pedidos-v1.0-entrega.zip}, que contiene el código, la configuración, los datos, los resultados, este informe y el notebook ejecutado.
El mismo contenido está disponible en el repositorio público \url{https://github.com/MauricioLopezC/io2-pedidos}, en la versión etiquetada \texttt{v1.0-entrega}.

\subsubsection*{Requisitos}

\begin{itemize}
  \item Python 3.12 o posterior; las versiones exactas de las bibliotecas están fijadas en \texttt{requirements.txt}.
  \item La reproducción se verificó con Python 3.14 en Linux.
  \item Para recompilar este informe, además, una distribución de \LaTeX{} con \texttt{latexmk} y \texttt{biber}.
\end{itemize}

\subsubsection*{Pasos}

Desde la carpeta \texttt{io2-pedidos}, se instalan las dependencias y el paquete del proyecto con \texttt{pip install -r requirements.txt}.
A continuación hay dos alternativas equivalentes:

\begin{enumerate}
  \item \textbf{Notebook.} Abrir \texttt{notebooks/recorrido.ipynb} con \texttt{jupyter notebook} y ejecutar todas las celdas.
  El notebook vuelve a correr el pipeline completo y compara la huella SHA-256 de cada archivo de resultados con la del archivo entregado.
  Después recorre los datos, el pronóstico, la verificación, los experimentos y la sensibilidad.
  \item \textbf{Línea de comandos.} Ejecutar \texttt{make all}, que corre los cinco scripts en orden y compila el informe.
\end{enumerate}

La ejecución completa demora del orden de dos minutos, la mayor parte en el análisis de sensibilidad.

\subsubsection*{Determinismo}

Toda la aleatoriedad usa generadores con semillas fijadas en \texttt{config/experimento.yaml}, y las semillas de las réplicas se derivan de ellas.
Por lo tanto, la misma configuración produce archivos de resultados y tablas idénticos byte a byte.
Las figuras PDF no se comparan de ese modo porque sus bytes dependen de la versión de Matplotlib, aunque los datos graficados sean los mismos.

La última sección del notebook permite cambiar el plazo, el cumplimiento exigido, la demanda y el tiempo de servicio sobre una copia de la configuración.
Esos resultados son exploratorios y no forman parte de este informe.

\subsubsection*{Origen de cada tabla y figura}

\begin{table}[htbp]
  \centering
  \caption{Script que genera cada tabla y figura del informe.}
  \label{tab:trazabilidad}
  \small
  \begin{tabular}{lll}
    \toprule
    Script & Comando & Tablas y figuras \\
    \midrule
    \texttt{01\_preparar\_datos.py} & \texttt{make datos} & Datos procesados y \texttt{resultados/datos\_resumen.csv} \\
    \texttt{02\_pronostico.py} & \texttt{make pronostico} & Tabla~\ref{tab:errores-pronostico} y Figura~\ref{fig:pronostico-prueba} \\
    \texttt{03\_verificacion.py} & \texttt{make verificacion} & Tablas~\ref{tab:servicio} y~\ref{tab:verificacion-mmc} \\
    \texttt{04\_experimentos.py} & \texttt{make experimentos} & Tablas~\ref{tab:experimentos-escenarios} y~\ref{tab:experimentos-recomendacion}, Figura~\ref{fig:experimentos-desempeno} \\
    \texttt{05\_sensibilidad.py} & \texttt{make sensibilidad} & Tabla~\ref{tab:sensibilidad-recomendacion} y Figura~\ref{fig:sensibilidad-cumplimiento} \\
    \bottomrule
  \end{tabular}
\end{table}

Cada script documenta en su encabezado los archivos que escribe en \texttt{resultados/}, \texttt{informe/tablas/} e \texttt{informe/figuras/}.
Los parámetros del estudio están todos en \texttt{config/experimento.yaml}, con un comentario que justifica cada valor.
